% The S6 record and the integral monodromy of its (3,4,infinity) period system (version 3).
% Prepared by Claude (Anthropic), model claude-opus-5-5 (Opus 5.5).
% Compile with lualatex (twice).
\documentclass[11pt,a4paper]{article}
\usepackage[a4paper,margin=24mm,headheight=14pt]{geometry}
\usepackage{fontspec}
\directlua{luaotfload.add_fallback("readerfallback",{"TeX Gyre Pagella Math:mode=harf;","STIX:mode=harf;","DejaVu Serif:mode=harf;","FreeSerif:mode=harf;"})}
\setmainfont{TeX Gyre Pagella}[RawFeature={fallback=readerfallback}]
\setsansfont{TeX Gyre Heros}[Scale=MatchLowercase]
\setmonofont{DejaVu Sans Mono}[Scale=0.84]
\usepackage{amsmath,amsthm,mathtools}
\usepackage{graphicx}
\usepackage{unicode-math}
\setmathfont{TeX Gyre Pagella Math}
\usepackage{microtype}
\usepackage{booktabs,longtable,array,tabularx}
\usepackage{enumitem}
\setlist{itemsep=2pt,topsep=4pt,leftmargin=*}
\usepackage[dvipsnames]{xcolor}
\usepackage{fancyhdr}
\usepackage{titlesec}
\usepackage[hidelinks,bookmarksnumbered]{hyperref}
\hypersetup{pdftitle={The S6 record and the integral monodromy of its (3,4,infinity) period system},
  pdfauthor={Claude (Anthropic), model claude-opus-5-5; record by KokunoYumeto},
  pdfsubject={The S6 verification and reconstruction record (Zenodo 10.5281/zenodo.22678442): results with proofs, and the integral monodromy of the period system},
  pdfkeywords={complex structure on the six-sphere, compact complex threefolds, algebraic dimension, torus fibrations, monodromy, paramodular group, Heisenberg group, arithmetic groups}}
\usepackage{bookmark}

\titleformat{\section}{\normalfont\Large\bfseries}{\thesection}{0.8em}{}
\titleformat{\subsection}{\normalfont\large\bfseries}{\thesubsection}{0.7em}{}
\titlespacing*{\section}{0pt}{2.2ex plus 1ex}{1.2ex}
\titlespacing*{\subsection}{0pt}{1.8ex plus .8ex}{0.9ex}

\pagestyle{fancy}
\fancyhf{}
\fancyhead[L]{\small\itshape The $S^6$ record and the integral monodromy of its period system}
\fancyhead[R]{\small\thepage}
\renewcommand{\headrulewidth}{0.3pt}
\fancypagestyle{plain}{\fancyhf{}\fancyfoot[C]{\small\thepage}\renewcommand{\headrulewidth}{0pt}}

\theoremstyle{plain}
\newtheorem{theorem}{Theorem}[section]
\newtheorem{proposition}[theorem]{Proposition}
\newtheorem{lemma}[theorem]{Lemma}
\newtheorem{corollary}[theorem]{Corollary}
\newtheorem{introthm}{Theorem}
\renewcommand{\theintrothm}{\Roman{introthm}}
\theoremstyle{definition}
\newtheorem{definition}[theorem]{Definition}
\newtheorem{example}[theorem]{Example}
\theoremstyle{remark}
\newtheorem{remark}[theorem]{Remark}

% status line printed under a statement
\newcommand{\status}[1]{\par\smallskip\noindent{\small\color{black!70}\textsf{Status.}\ #1}\par\smallskip}
\newcommand{\negres}[2]{\par\smallskip\noindent\textbf{#1}\ #2\par}

% notation
\newcommand{\B}{\mathcal{B}}
\newcommand{\Zs}{\mathcal{Z}}
\newcommand{\C}{\mathbb{C}}
\newcommand{\R}{\mathbb{R}}
\newcommand{\Q}{\mathbb{Q}}
\newcommand{\ZZ}{\mathbb{Z}}
\newcommand{\NN}{\mathbb{N}}
\newcommand{\Ff}{\mathbb{F}}
\newcommand{\Pp}{\mathbb{P}}
\newcommand{\re}{\operatorname{Re}}
\newcommand{\im}{\operatorname{Im}}
\newcommand{\Res}{\operatorname*{Res}}
\newcommand{\sgn}{\operatorname{sgn}}
\newcommand{\Gr}{\operatorname{Gr}}
\newcommand{\cl}{\overline}
\newcommand{\brkus}{\textunderscore\allowbreak}
\newcommand{\note}[1]{{\ttfamily\let\_\brkus #1}}
\newcommand{\fn}[1]{{\ttfamily\let\_\brkus #1}}
\newcommand{\lbl}[1]{\textsf{#1}}

\setlength{\parskip}{3pt}
\setlength{\emergencystretch}{2.5em}
\setlength{\parindent}{0pt}

\newcommand{\T}{\mathcal{T}}
\newcommand{\Hs}{\mathcal{H}}
\newcommand{\lcm}{\operatorname{lcm}}
\newcommand{\TV}{\operatorname{TV}}
\newcommand{\tr}{\operatorname{tr}}
\newcommand{\Var}{\operatorname{Var}}
